\documentclass[11pt]{article}
\usepackage[margin=1in]{geometry}
\usepackage{graphicx}
\usepackage{setspace}
\usepackage{enumitem}
\usepackage{amsmath}
\usepackage[T1]{fontenc}
\usepackage{hyperref}

\title{Check-In \#2\\
\large Fate of the Solar Eclipse: Evolution of the Moon-Sun Angular Size Ratio}
\author{Clover Lee}
\date{Spring 2026 \\ \small Due Friday, April 10th, 2026}

\begin{document}
\maketitle

% ─────────────────────────────────────────────────────────────
\section*{1. Feedback \& How You Applied It}
% What feedback did you receive on your proposal, and how did
% you apply it to your project?
% ─────────────────────────────────────────────────────────────

[Summarize the feedback you received on your proposal here.
Then describe 1--2 specific changes you made in response.
For example: did a TA suggest narrowing the scope, improving
the data cleaning step, or adding a particular visual?]


% ─────────────────────────────────────────────────────────────
\section*{2. Progress: Completed \& Still To Do}
% What have you completed so far? What do you still need to finish?
% ─────────────────────────────────────────────────────────────

\textbf{Completed so far:}
\begin{itemize}
    \item Defined physical constants and the linear recession model
          $d_\mathrm{Moon}(t) = d_0 + vt$
    \item Set up Jupyter notebook structure and imports
    \item [Add any other completed items here]
\end{itemize}

\textbf{Still to finish:}
\begin{itemize}
    \item Successfully load and clean the NASA Five Millennium eclipse catalog
    \item Generate Plot 1 (Moon angular diameter vs.\ time)
    \item Generate Plot 2 ($R(t)$ ratio with threshold and shading)
    \item Generate Plot 3 (superimposed Sun/Moon disk visualization)
    \item Add \texttt{ipywidgets} interactive slider
    \item Write final report and polish notebook
\end{itemize}


% ─────────────────────────────────────────────────────────────
\section*{3. How the Project Has Changed}
% In what ways has your project changed since your original proposal?
% ─────────────────────────────────────────────────────────────

[Describe any changes in scope, methods, or focus since your
proposal. For example: Did the NASA catalog turn out to be
harder to parse than expected? Did you add or drop a planned
visualization? Did you decide to use or not use \texttt{astropy}?]


% ─────────────────────────────────────────────────────────────
\section*{4. Most Challenging Part}
% What has been the most challenging part of the project?
% ─────────────────────────────────────────────────────────────

[Describe the single biggest challenge you have encountered.
This could be technical (e.g., parsing the fixed-width NASA file,
getting the disk visualization to scale correctly) or conceptual
(e.g., deciding on a time range, understanding the $\Delta T$
correction). Be specific.]


% ─────────────────────────────────────────────────────────────
\section*{5. New Things Learned}
% What are 2–3 new things you've learned while working on this project?
% ─────────────────────────────────────────────────────────────

\begin{enumerate}
    \item \textbf{[Topic 1]:} [Description — e.g., how \texttt{pd.read\_fwf}
          handles fixed-width files, or what the Saros cycle is]
    \item \textbf{[Topic 2]:} [Description — e.g., how \texttt{matplotlib.patches.Circle}
          works, or what eclipse magnitude actually measures]
    \item \textbf{[Topic 3 (optional)]:} [Description]
\end{enumerate}


% ─────────────────────────────────────────────────────────────
%  ROUGH DRAFT REPORT
% ─────────────────────────────────────────────────────────────
\newpage
\begin{center}
    {\Large \textbf{Rough Draft Report}}\\[0.5em]
    {\large Fate of the Solar Eclipse: Evolution of the Moon-Sun Angular Size Ratio}\\[0.3em]
    Clover Lee \quad Spring 2026
\end{center}

\section*{Introduction}

A solar eclipse occurs when the Moon passes between the Earth and the Sun.
Total solar eclipses — in which the Moon completely obscures the Sun's disk —
are only possible when the Moon's apparent angular size is at least as large as the Sun's.
The angular diameter of an object is approximated by the small-angle formula:
\[
    \theta \approx \frac{D}{d}
\]
where $D$ is the object's true diameter and $d$ is its distance from the observer.

The Moon is gradually receding from Earth due to tidal interactions at a rate of
approximately $3.8\,\text{cm/year}$, causing its apparent angular size to decrease
slowly over time.
This raises a natural question: \textit{how long will total solar eclipses remain
possible, and when will only annular eclipses occur?}

\section*{Background}

[Expand on the physics here in your own words.
You might discuss:
\begin{itemize}
    \item Why the Moon recedes (tidal friction / angular momentum transfer)
    \item The difference between total, annular, and hybrid eclipses
    \item Why this calculation uses a linear model (and its limitations)
\end{itemize}
This section can grow as you learn more.]

\section*{Methods}

\subsection*{Model}

The Moon's distance from Earth is modeled as a linear function of time:
\[
    d_\mathrm{Moon}(t) = d_0 + v \cdot (t - t_0)
\]
where $d_0 = 384{,}400\,\text{km}$ is the current mean Earth-Moon distance,
$v = 3.8 \times 10^{-5}\,\text{km/yr}$ is the recession rate,
and $t_0 = 2026$.

The angular size ratio is:
\[
    R(t) = \frac{\theta_\mathrm{Moon}(t)}{\theta_\mathrm{Sun}}
         = \frac{D_\mathrm{Moon} / d_\mathrm{Moon}(t)}{D_\mathrm{Sun} / d_\mathrm{Sun}}
\]

Total eclipses are possible when $R(t) \geq 1$, and only annular eclipses
occur when $R(t) < 1$.
Setting $R = 1$ and solving analytically gives the crossover time:
\[
    t^* = \frac{D_\mathrm{Moon} \cdot d_\mathrm{Sun} / D_\mathrm{Sun} - d_0}{v}
\]

\subsection*{Data}

Historical eclipse data was obtained from NASA's \textit{Five Millennium Canon of
Solar Eclipses} \cite{nasa_canon}, a fixed-width text catalog spanning from
$-1999$ to $+3000$.
Key fields used from the catalog include \texttt{Year}, \texttt{Ecl\_Type}
(eclipse type: T/A/H/P), and \texttt{Ecl\_Mag} (eclipse magnitude,
equivalent to $R$ at the moment of each eclipse).
The catalog was loaded with \texttt{pandas.read\_fwf} with \texttt{latin1}
encoding and cleaned by converting relevant columns to numeric types.

\section*{Results}

[This section will be filled in once your plots are complete.
For now, describe what you have so far:]

The analytic crossover occurs at approximately $t^* \approx$ [your computed value]
million years from today, corresponding to the year [year].
At the present day ($t = 0$), the model gives $R(2026) \approx$ [your computed value],
consistent with the fact that total eclipses are currently still possible.

\textbf{[Insert Plot 1 here: Moon angular diameter vs.\ time]}

\textbf{[Insert Plot 2 here: $R(t)$ ratio with threshold]}

\textbf{[Insert Plot 3 here: superimposed disk comparison]}

\section*{Discussion}

[To be expanded. Address questions such as:
\begin{itemize}
    \item Does your computed crossover year match published estimates
          ($\sim$600 million years)?
    \item What are the main sources of error or simplification
          (e.g., linear recession, ignoring orbital eccentricity,
          ignoring Earth-Sun distance variation)?
    \item How does the NASA eclipse magnitude data compare to your model's
          $R(t)$ curve?
    \item Is the linear recession assumption physically justified over
          such long timescales?
\end{itemize}]

\section*{Conclusion}

[To be written after results are finalized.
One or two sentences summarizing the key finding:
when will the last total solar eclipse occur, and what does
this mean physically?]

\begin{thebibliography}{9}
\bibitem{nasa_canon}
    NASA Goddard Space Flight Center.
    \textit{Five Millennium Canon of Solar Eclipses: -1999 to +3000}.
    \url{https://eclipse.gsfc.nasa.gov/5MCSE/5MKSEcatalog.txt}

\bibitem{chapront}
    Chapront, J., Chapront-Touze, M., \& Francou, G. (2002).
    A new determination of lunar orbital parameters, precession constant,
    and tidal acceleration from LLR measurements.
    \textit{Astronomy \& Astrophysics}, 387, 700--709.

\bibitem{astropy}
    Astropy Collaboration (2022).
    The Astropy Project.
    \url{https://www.astropy.org}
\end{thebibliography}

\end{document}
